\documentclass[a4paper,12pt]{article}

% Encodage et langue
\usepackage[utf8]{inputenc}
\usepackage[T1]{fontenc}
\usepackage[french]{babel}

% Graphiques et mise en page
\usepackage{graphicx}
\usepackage{geometry}
\usepackage{array}
\usepackage{float}
\usepackage{xcolor}
\usepackage{titlesec}
\usepackage{setspace}
\usepackage{hyperref}
\geometry{margin=2.5cm}
\setstretch{1.15}
\hypersetup{colorlinks=true, linkcolor=black, urlcolor=blue}

\newcommand{\sectionbreak}{\clearpage}
\newcommand{\newsubsection}[1]{\clearpage\subsection{#1}}
\newcommand{\imageplaceholder}[2]{%
    \begin{figure}[H]
        \centering
        \fbox{%
            \parbox[c][3.8cm][c]{0.88\textwidth}{%
                \centering
                \textbf{Emplacement d'image à remplacer}\\[0.35cm]
                \emph{#2}
            }%
        }
        \caption{#1}
    \end{figure}
}

\setcounter{tocdepth}{2}

\title{\textbf{Chamouraï} \\ Rapport de soutenance}
\author{Les Chatpitres}
\date{}

\begin{document}

\maketitle

\begin{center}
    \includegraphics[width=0.6\textwidth]{titre.png}

    \vspace{1cm}

    \includegraphics[width=0.4\textwidth]{logo.png}
\end{center}

\newpage
\tableofcontents
\newpage

\section{Introduction}

\textit{Chamouraï} est un jeu d'action et d'aventure en vue du dessus, réalisé en Python avec
Pygame. Le joueur incarne un chat samouraï qui doit progresser d'une forêt vers un
donjon, affronter des ennemis de plus en plus dangereux, réunir deux clés puis vaincre
un golem. Notre objectif était de proposer une expérience immédiatement jouable, mais
qui demande de maîtriser progressivement le placement, l'esquive et le rythme des
attaques.

Le projet a beaucoup évolué depuis sa première version. Le prototype initial reposait
principalement sur le déplacement du joueur, l'épée orientée à la souris, quelques
ennemis et une carte. La version actuelle comprend un apprentissage guidé, des objectifs
successifs, plusieurs familles d'ennemis, des adversaires majeurs, un verrouillage du donjon par
clés, une carte consultable, des musiques contextuelles et un mode à deux joueurs avec
gestion de la mort et de la reprise de partie.

Le jeu repose sur un principe d'apprentissage par l'échec : perdre fait partie de
l'apprentissage. Nous avons cependant cherché à rendre les échecs compréhensibles.
Les animations annoncent les attaques, la barre de vie de l'adversaire majeur donne un repère clair
et la reprise depuis le début du donjon évite de répéter inutilement l'introduction.

L'identité de \textit{Chamouraï} vient du rapprochement de deux univers que nous
apprécions : l'ambiance de samouraïs et les paysages japonais de
\textit{Ghost of Tsushima}, associés au personnage félin de \textit{Stray}. Cette
association nous a conduit à imaginer un héros à la fois attachant et combattant :
un chat samouraï placé dans un monde menacé, mais encore riche en couleurs et en
détails.

\imageplaceholder{Vue générale d'une partie de \textit{Chamouraï}}{Capture du joueur dans la forêt, avec plusieurs squelettes visibles, l'interface de vie et le panneau d'objectif en haut à gauche.}

\sectionbreak
\section{Répartition des rôles}

\begin{center}
\begin{tabular}{|m{4.2cm}|m{9cm}|}
\hline
\textbf{Membre} & \textbf{Responsabilités principales} \\
\hline
Baptiste LOISEAU & Chef de projet, conception des cartes, jouabilité, intégration graphique et équilibrage des combats \\
\hline
Lucas HUET & Interfaces, menus, musiques, effets sonores et réglages du son \\
\hline
Loan BLANPAIN & Site internet, scénario, univers et communication du projet \\
\hline
Simon BOURGEOIS & Intelligence artificielle, déplacements des ennemis et multijoueur \\
\hline
\end{tabular}
\end{center}

Cette répartition nous a donné des domaines de responsabilité clairs, mais le jeu ne
peut pas être assemblé en isolant complètement les tâches. Par exemple, une nouvelle
zone créée dans Tiled influence les collisions, l'intelligence artificielle, les
objectifs, les sons et la synchronisation réseau. Une partie importante du travail a
donc été l'intégration : vérifier ensemble les ajouts de chacun et résoudre les conflits
lorsque plusieurs fonctionnalités modifiaient la boucle principale du jeu.

\sectionbreak
\section{Parcours de jeu et fonctionnalités}

\subsection{Une progression guidée}

Au lancement d'une partie, le joueur arrive dans la zone extérieure. Plutôt que de
lui demander de découvrir toutes les commandes au milieu d'un combat difficile, nous
avons ajouté une suite d'objectifs affichée en haut à gauche de l'écran :

\begin{enumerate}
    \item parler au vieil homme afin de découvrir le parchemin et l'interaction ;
    \item frapper trois fois afin de comprendre l'utilisation de l'épée ;
    \item esquiver trois fois pour apprendre la mobilité défensive ;
    \item éliminer trois squelettes dans la forêt ;
    \item entrer dans le donjon ;
    \item ouvrir la carte avec la touche \texttt{M} ;
    \item trouver les deux clés permettant d'ouvrir la grille.
\end{enumerate}

Le portail vers le donjon n'est accessible que lorsque les premiers objectifs sont
validés. Ce choix évite qu'un joueur se retrouve face aux adversaires majeurs sans avoir utilisé les
actions essentielles. À deux joueurs, les éliminations utiles à l'apprentissage sont
synchronisées : les deux joueurs partagent ainsi la même progression au lieu de se
retrouver bloqués à des étapes différentes.

\imageplaceholder{Apprentissage et objectifs successifs}{Capture du début de partie avec le panneau d'objectif affichant « Parler au vieil homme », puis une étape chiffrée comme « Frapper 3 fois ».}

\subsection{Exploration et carte}

Dans le niveau principal, la touche \texttt{M} affiche une carte simplifiée. Elle
permet au joueur de localiser sa position dans un environnement plus vaste, sans
supprimer la nécessité d'explorer. La carte est volontairement fonctionnelle :
elle montre les zones utiles et le joueur, mais ne remplace pas le déplacement réel
dans le niveau.

La grille qui ferme l'accès au golem constitue l'objectif central du donjon. Deux adversaires majeurs
intermédiaires laissent tomber chacun une clé lorsqu'ils sont vaincus. Dans la zone
devant la grille, un message indique le nombre de clés requis ; après collecte des
deux clés, le joueur peut appuyer sur \texttt{E} pour ouvrir le passage. Cette
interaction relie combat et exploration : éliminer les adversaires majeurs a une conséquence visible
sur la carte.

\imageplaceholder{Carte du niveau et position du joueur}{Capture de la carte ouverte avec la position du joueur, les principales salles et la zone de l'adversaire final.}

\imageplaceholder{Grille et compteur de clés}{Capture devant la grille fermée, avec les deux emplacements de clés à l'écran et le message demandant deux clés.}

\subsection{Narration environnementale}

Un vieil homme animé est présent dans la zone de départ. Lorsque le joueur s'approche,
une icône animée de touche \texttt{E} apparaît ; l'interaction ouvre un parchemin en
grand à l'écran. Pendant sa lecture, le déplacement, l'esquive et l'attaque sont
bloqués afin de ne pas provoquer une action involontaire derrière l'interface.

Le parchemin a deux rôles. Il présente d'abord les commandes principales au joueur,
puis il introduit la situation : la population est terrorisée par les monstres du
donjon, qui commencent désormais à rôder jusque dans la forêt. Cette scène donne ainsi
un objectif au personnage avant son premier affrontement, sans imposer une longue
séquence d'explication séparée du jeu. Nous avons ralenti l'animation de la touche
d'interaction afin qu'elle attire l'attention sans clignoter de manière gênante.

\imageplaceholder{Interaction avec le vieil homme}{Capture du vieil homme avec l'icône E visible, puis du parchemin ouvert en superposition à l'écran.}

\sectionbreak
\section{Conception technique des cartes}

\subsection{Cartes produites avec Tiled}

Les niveaux sont dessinés avec Tiled, qui permet de superposer des calques de tuiles et
des calques d'objets. Le programme charge les fichiers \texttt{.tmx} et convertit leur
contenu en éléments de jeu. Cette méthode est beaucoup plus souple que d'écrire les
positions directement en Python : la carte peut être retouchée visuellement sans
modifier la logique centrale.

La forêt devait être davantage qu'une simple zone de départ. Nous souhaitions créer
un environnement immersif et particulièrement détaillé, afin que le joueur ressente
la différence entre un extérieur encore vivant, mais désormais dangereux, et
l'intérieur minéral du donjon. Le travail par calques dans Tiled nous a permis
d'ajouter du relief visuel et des éléments de décor sans perdre la lisibilité des
chemins et des combats.

Pour les ennemis, chaque objet Tiled portant le nom \texttt{Skeleton},
\texttt{Necromancer}, \texttt{SkeletonBoss}, \texttt{NecromancerBoss} ou
\texttt{Golem} crée automatiquement le monstre correspondant. Ajouter un squelette
revient donc à déposer un nouvel objet sur la carte. Le même principe est utilisé pour
les zones d'apparition et d'interaction.

\imageplaceholder{Édition du niveau dans Tiled}{Capture de Tiled montrant les objets Skeleton, NecromancerBoss et Golem placés sur le niveau ainsi que les calques Walls, Grid et GridKey.}

\subsection{Collisions des murs et de la grille}

La première carte utilisait des collisions définies objet par objet, ce qui devenait
long et source d'erreurs dès que le décor changeait. Nous avons choisi de rendre
bloquante toute tuile non vide du calque \texttt{Walls}. Le responsable de la conception des niveaux peut ainsi
dessiner un mur et obtenir immédiatement une collision cohérente.

L'accès à l'adversaire final a nécessité une distinction importante. Le calque de tuiles
\texttt{Grid} contient l'apparence et les collisions de la grille, tandis que le
calque d'objets \texttt{GridKey} définit uniquement la zone où le message
d'interaction apparaît. Confondre ces deux types de calques entraînait une erreur au
chargement, car une tuile fournit des coordonnées et un identifiant tandis qu'un
objet Tiled fournit un rectangle. Après ouverture, le calque visuel de la grille est
masqué et ses collisions sont retirées simultanément.

\subsection{Pourquoi ce choix est durable}

La génération pilotée par la carte réduit les dépendances entre la programmation et la
création des niveaux. Elle impose cependant une convention : les noms des calques et
des objets doivent être exacts. Nous avons accepté cette contrainte, car elle est
facile à documenter et bien plus pratique que d'ajouter une ligne de code pour chaque
nouvel ennemi ou obstacle.

\sectionbreak
\section{Système de combat}

\subsection{Épée orientée par la souris}

Le combat repose sur une coordination clavier-souris. Le joueur se déplace au clavier
et indique avec le curseur l'endroit vers lequel il souhaite orienter son épée. Pour
obtenir cette orientation, nous considérons le segment qui relie le personnage au
curseur. Ce segment forme un angle avec l'axe horizontal de l'écran.

La méthode utilisée est trigonométrique : à partir de l'écart horizontal et de l'écart
vertical entre le joueur et la souris, une fonction d'arc tangente permet de retrouver
l'angle de visée. Cette méthode a l'avantage de fonctionner dans toutes les directions,
y compris en diagonale, contrairement à un système limité à « gauche » et « droite ».
L'épée peut ainsi suivre naturellement le curseur autour du personnage.

Une difficulté est apparue lorsque le joueur changeait très rapidement de côté pendant
une attaque : l'animation pouvait sembler recommencer. Nous avons donc choisi de
conserver l'orientation prise au début d'un coup jusqu'à sa fin. Pour le joueur, le
résultat est plus clair : chaque action déclenche une seule animation complète.

Cette mécanique a constitué l'une des parties les plus difficiles du projet. À la
différence d'une attaque dirigée dans seulement quatre sens, elle demande de lier en
continu le mouvement de la souris à l'orientation de l'arme. Elle nous a notamment
obligés à mobiliser des notions de géométrie et de trigonométrie pour obtenir un
comportement fluide et compréhensible.

\imageplaceholder{Attaque orientée par le curseur}{Deux captures côte à côte : le joueur attaque vers la droite puis vers la gauche, avec le curseur visible.}

\subsection{Esquive directionnelle}

L'esquive fonctionnait initialement seulement sur un axe horizontal : regarder à
droite envoyait le joueur vers la gauche et inversement. Ce comportement était
insuffisant pour les combats en arène. L'esquive utilise désormais le vecteur opposé
à la souris : un curseur placé en bas à droite provoque une esquive vers le haut à
gauche.

Pendant l'esquive, les déplacements normaux sont bloqués. Sans ce verrouillage, le
joueur pouvait modifier artificiellement sa trajectoire et rendre l'action difficile
à anticiper. Cette décision améliore aussi l'équilibrage : l'esquive est une action
forte, mais elle engage le joueur pendant une courte durée.

\sectionbreak
\section{Ennemis et adversaires majeurs}

\subsection{Squelettes et nécromanciens}

Le squelette constitue l'ennemi rapproché de base : il poursuit le joueur et attaque à
courte portée. Le nécromancien introduit une menace à distance et oblige à se déplacer
plutôt qu'à attendre derrière son épée. Chacun existe également en version renforcée, avec
une apparence agrandie, une barre de vie visible au-dessus de sa tête, davantage de vie,
une vitesse supérieure et des attaques plus fréquentes.

Lorsque plusieurs squelettes poursuivaient le joueur, ils avaient tendance à se
superposer. Un seul coup pouvait alors toucher tout le groupe sans réel risque. Nous
avons ajouté des zones de collision entre ennemis et un espacement dans leur recherche
de trajectoire. L'épée peut toujours toucher plusieurs ennemis si son arc les atteint,
ce qui récompense un bon placement, mais le groupe n'arrive plus comme une masse
immobile.

\subsection{Recherche de chemin}

Les monstres utilisent une recherche de chemin basée sur A*. Nous avons préféré A* à
Dijkstra car notre objectif n'est pas d'explorer toutes les distances possibles, mais
de diriger efficacement un ennemi vers une cible connue en tenant compte des murs. Son
heuristique limite la quantité de cases inspectées et convient à une carte découpée en
tuiles.

L'intégration a été délicate : un ennemi doit éviter les murs de Tiled, garder une
distance raisonnable avec les autres ennemis et continuer à attaquer lorsque le joueur
est accessible. Un trajet mathématiquement court n'est pas toujours agréable en jeu,
notamment lorsque tous les ennemis choisissent exactement le même passage. L'espacement
ajouté aux squelettes répond donc à un besoin de jouabilité autant qu'à un problème de
collision.

\imageplaceholder{Affrontement avec différents ennemis}{Capture d'une salle montrant des squelettes espacés et un nécromancien attaquant à distance.}

\subsection{Adversaires intermédiaires et clés}

Le \textit{SkeletonBoss} et le \textit{NecromancerBoss} ne sont pas seulement des
ennemis agrandis : leurs statistiques renforcées donnent deux combats intermédiaires
qui préparent à l'adversaire final. À leur mort, ils déposent chacun une clé représentée par
un objet ramassable. Ce choix rend la progression immédiatement lisible : le joueur
voit l'objet obtenu et le compteur de clés confirme son avancement.

\imageplaceholder{Adversaire intermédiaire et clé déposée}{Capture d'un adversaire majeur avec sa barre de vie, puis de la clé apparue au sol après sa défaite.}

\subsection{Le golem, adversaire final}

Le golem est le point culminant du niveau. Son échelle est volontairement imposante et
sa barre de vie est affichée en haut de l'écran dès que la grille tombe. En première
phase, il dispose d'une attaque de mêlée et d'un tir de bras. À partir de la moitié de
sa vie, son apparence devient lumineuse, ses statistiques augmentent et il obtient une
attaque supplémentaire : une onde de choc bleue en arc de cercle.

Le choix du golem est aussi issu d'une contrainte de production. Nous avions décidé
d'utiliser des ressources graphiques gratuites ; parmi celles qui correspondaient à
l'esthétique du jeu, le golem était le seul adversaire disposant d'un ensemble
d'attaques assez complet et varié pour construire un véritable combat final. Cette
contrainte est devenue une opportunité : sa silhouette imposante correspond bien à
une menace enfermée au fond du donjon.

Nous avions envisagé un rayon lumineux continu sortant de sa tête. Sa gestion s'est
révélée difficile à rendre fiable et lisible : placement exact de l'origine, feuille
d'animation disposée verticalement, collision d'un rayon long et synchronisation des
images d'animation. Nous avons remplacé cette idée par une onde de choc. Elle conserve le rôle
recherché, c'est-à-dire forcer le joueur à esquiver en seconde phase, mais sa forme,
sa vitesse et sa zone de contact sont plus simples à comprendre et à régler.

La mort du golem joue une animation plus longue et laisse son corps visible, ce qui
donne du poids à la victoire. L'écran de fin propose ensuite un retour au menu
principal.

\imageplaceholder{Combat contre le golem en phase 2}{Capture du golem lumineux, avec sa barre de vie en haut de l'écran et son onde de choc bleue en arc de cercle.}

\imageplaceholder{Écran de victoire}{Capture de l'écran affiché après la mort du golem, avec le bouton permettant de retourner au menu principal.}

\sectionbreak
\section{Interface, graphismes et son}

\subsection{Direction graphique et animations}

Le jeu adopte un style de dessin en pixels en vue du dessus. Le personnage principal a été
généré puis retravaillé pour correspondre à l'univers ; ses animations ont été
adaptées aux actions principales. Les feuilles de tuiles et plusieurs monstres ont
été sélectionnés dans des ressources compatibles avec notre direction artistique,
puis intégrés à nos contraintes d'échelle et d'animation.

\begin{figure}[H]
    \centering
    \includegraphics[width=0.6\textwidth]{animation_joueur.png}
    \caption{Feuille d'animation de marche du joueur}
\end{figure}

L'ajout du vieil homme, des clés et des différents adversaires majeurs a demandé une attention
particulière aux dimensions des images : un adversaire majeur doit sembler plus massif tout en
conservant des collisions jouables. Le golem, notamment, utilise des images plus
grandes que les autres personnages et des images séparées pour ses projectiles.

\subsection{Retours des personnes ayant essayé le jeu}

Les personnes qui ont essayé \textit{Chamouraï} ont particulièrement apprécié
l'association entre le chat et l'univers samouraï, qui donne rapidement une identité
au projet. Le système de combat guidé par la souris a également été perçu comme
original. Ces retours ont confirmé l'intérêt de conserver une attaque libre et
directionnelle, même si sa réalisation était plus complexe qu'un système de combat
orienté uniquement dans quelques directions fixes.

\subsection{Menus et informations en jeu}

Le menu principal permet d'accéder au jeu à un joueur, au jeu à deux, aux options et à la
sortie. Le menu du jeu à deux offre la création ou la connexion à une partie et comprend
un bouton de retour, indispensable pour ne pas enfermer le joueur dans un parcours de
connexion. En jeu, l'interface affiche les informations utiles au bon moment :
objectifs pendant l'apprentissage, clés dans le donjon, carte sur demande et barre du
golem lors du combat final.

\begin{figure}[H]
    \centering
    \includegraphics[width=0.6\textwidth]{interface.png}
    \caption{Interface principale du jeu}
\end{figure}

\imageplaceholder{Menu du jeu à deux}{Capture du menu permettant de rejoindre ou créer une partie, avec le bouton Retour visible.}

\subsection{Musiques et effets sonores}

Les sons et musiques renforcent la compréhension de l'action : bruit d'attaque,
ambiances et transitions musicales accompagnent les différents états de partie. La
musique du donjon est lancée lors de l'entrée dans le niveau ou d'une reprise, puis une
musique de combat final prend le relais lorsque la grille s'ouvre. Le menu des sons permet de
régler l'expérience sans modifier le code.

Nous avons également différencié les sons selon le lieu : les pas sur l'herbe dans
la zone extérieure ne sont pas ceux des pas sur la pierre du donjon. Les musiques
suivent la même progression dramatique : une ambiance calme accompagne la forêt,
une musique plus dynamique souligne l'entrée dans le donjon, puis une musique épique
marque l'affrontement contre le golem.

L'intégration sonore a posé un problème fréquent dans un projet en équipe : la boucle
principale était également modifiée par les objectifs, les adversaires majeurs et le jeu à deux.
Lors de la fusion de ces travaux, il fallait préserver les appels sonores sans perdre
les règles de progression ou de reprise. Nous avons donc rattaché la musique à des
événements identifiables, comme une transition de carte ou l'ouverture de la grille,
plutôt que de déclencher un son de manière diffuse à chaque rafraîchissement d'écran.

\begin{figure}[H]
    \centering
    \includegraphics[width=0.6\textwidth]{sons.png}
    \caption{Organisation des ressources sonores}
\end{figure}

\sectionbreak
\section{Mode à deux joueurs en coopération}

\subsection{Principe de connexion}

Le mode à deux joueurs est conçu pour la coopération. L'un héberge la partie et l'autre rejoint
la session ; l'hôte exécute la logique serveur nécessaire à la cohérence de la partie.
Les joueurs envoient leurs actions et leurs états, tandis que les éléments qui doivent
être identiques pour tous, notamment les monstres, leurs points de vie, les clés et
l'ouverture de la grille, sont synchronisés depuis le serveur.

Ce choix est adapté à un projet de petite taille : il évite une infrastructure en ligne
permanente tout en offrant une expérience en coopération. Il demande cependant de ne pas
faire confiance uniquement à l'affichage local. Si chaque ordinateur déplaçait ses
propres monstres, un adversaire majeur pourrait être mort sur un écran et vivant sur l'autre.

Nous avons volontairement limité ce mode à deux participants. Un mode pour quatre
joueurs aurait aussi exigé de prendre en charge les parties à trois et d'équilibrer
chaque configuration séparément. Nos cartes sont conçues pour un ou deux personnages :
au-delà, l'espace de combat deviendrait encombré et les animations se superposeraient,
ce qui rendrait l'action plus difficile à lire. Enfin, l'ordinateur de l'hôte gère le
serveur et les échanges de la partie ; multiplier les participants aurait accru sa
charge de traitement et la complexité de la synchronisation.

\subsection{Progression et équilibrage à deux}

La téléportation de la zone extérieure vers le donjon exige désormais que les deux
joueurs soient présents dans la zone. Lorsqu'un seul joueur attend, un message
l'indique clairement. Cette règle empêche un joueur de changer brutalement la carte
de son partenaire pendant qu'il termine un combat ou lit un parchemin.

À deux, tous les monstres possèdent deux fois plus de vie. Sans cette adaptation, les
deux épées réduiraient fortement la durée des affrontements, en particulier contre les
adversaires majeurs. Le doublement des points de vie est un réglage simple, cohérent et compréhensible
pour maintenir l'intérêt des combats en coopération.

\subsection{Mort, observation et reprise}

Lorsqu'un joueur meurt dans une partie à deux, il disparaît du combat et peut observer la
partie : sa caméra suit le partenaire encore vivant. Il peut aussi retourner au menu.
Si les deux joueurs meurent, un écran de défaite apparaît. Chacun peut voter pour
recommencer ; la partie ne redémarre au début du donjon que lorsque les deux votes ont
été reçus.

Cette fonction a nécessité de synchroniser davantage qu'une simple position : état
vivant ou mort, vote, remise à zéro des ennemis, des clés et de la grille, point de
réapparition et musique. Nous avons retenu un redémarrage au début du donjon plutôt
qu'au tout début du jeu afin de respecter le principe d'apprentissage par l'échec sans
imposer à nouveau l'apprentissage guidé.

\imageplaceholder{Partie en coopération et joueur observateur}{Capture d'une partie à deux où un joueur combat tandis que l'autre, mort, observe la partie avec sa caméra centrée sur son partenaire.}

\imageplaceholder{Défaite et vote de reprise}{Capture de l'écran de défaite à deux joueurs présentant le bouton de retour et le bouton de vote pour recommencer.}

\sectionbreak
\section{Difficultés rencontrées et choix de conception}

\subsection{Faire évoluer une carte sans réécrire le jeu}

Le placement en dur des monstres était rapide pour un prototype, mais devenait un
frein dès que la carte changeait. L'utilisation des objets Tiled a demandé de définir
des conventions de nommage et de gérer plusieurs types de monstres, mais elle permet
aujourd'hui d'ajouter une rencontre sans toucher au code. De même, les collisions par
calque ont transformé une tâche manuelle et fragile en une opération de conception de niveau.

\subsection{Concilier lisibilité et ambition des adversaires majeurs}

Les attaques spectaculaires sont intéressantes seulement si le joueur comprend comment
y répondre. Le rayon lumineux du golem était ambitieux, mais son rendu et sa collision
manquaient de fiabilité. Le passage à une onde de choc nous a appris qu'une mécanique
plus maîtrisée est préférable à un effet complexe qui semble injuste. La seconde phase
reste marquante grâce au changement visuel, aux statistiques renforcées et à la
nouvelle attaque.

\subsection{Synchroniser des systèmes devenus interdépendants}

L'arrivée du jeu à deux a révélé que de nombreux systèmes étaient liés : vaincre un adversaire majeur
produit une clé, deux clés ouvrent une grille, l'ouverture lance la musique et permet
d'affronter le golem, puis la mort du golem affiche la victoire. En réseau, toutes ces
transitions doivent être partagées dans le même ordre. Nous avons donc déplacé vers
l'hôte les décisions importantes et transmis les états nécessaires aux postes connectés.

Les reprises de partie ont accentué cette difficulté. Relancer uniquement les joueurs
laisserait une grille déjà ouverte ou un adversaire majeur déjà vaincu. Le redémarrage doit remettre
en cohérence tout l'état du donjon. Ce travail est moins visible qu'une image animée ou un
menu, mais il conditionne la fiabilité de l'expérience en coopération.

\subsection{Intégrer le travail de plusieurs branches}

Le développement en parallèle a permis d'avancer rapidement sur le son, les ennemis,
l'apprentissage guidé et le réseau. En contrepartie, certaines fusions ont créé des conflits
dans le fichier principal, car toutes ces fonctionnalités interviennent dans la boucle
du jeu. La résolution ne consistait pas uniquement à choisir une version : il fallait
recomposer les comportements pour conserver à la fois les musiques, les objectifs,
les boutons de reprise et les nouvelles règles du jeu à deux.

Cette expérience justifie notre effort de séparation progressive des responsabilités :
cartes et collisions dans la gestion de niveau, ennemis dans leurs classes, menu et
son dans l'interface, synchronisation dans le serveur. Plus chaque fonction possède
un emplacement clair, plus une future évolution est vérifiable et intégrable.

\sectionbreak
\section{Communication et organisation}

\subsection{Outils de communication}

Le serveur \textit{Discord} centralise les échanges techniques, les ressources et les
réunions vocales. Il permet de signaler rapidement une anomalie, de montrer une capture ou
de demander une vérification après une intégration. Instagram est utilisé de manière
complémentaire pour des messages rapides et les imprévus.

\begin{figure}[H]
    \centering
    \includegraphics[width=0.58\textwidth]{discord.png}
    \caption{Serveur Discord du groupe Les Chatpitres}
\end{figure}

\subsection{Gestion des versions}

Nous utilisons Git pour répartir le développement et conserver l'historique du
projet. Afin de réduire les conflits, nous communiquons sur Discord avant d'intervenir
sur une même branche ou sur un même fichier important. Les images et les sons étant
volumineux, nous utilisons également Git LFS, une extension de Git adaptée au suivi
de ces fichiers.

Malgré cette organisation, les fusions de branches ont parfois été délicates. Le
fichier principal rassemble plusieurs comportements du jeu, comme les sons, les
objectifs ou le jeu à deux ; deux ajouts indépendants peuvent donc modifier les mêmes
passages. Certains conflits ont demandé de recomposer manuellement les modifications,
avec le risque de perdre une fonction déjà réalisée. Cette difficulté nous a appris à
mieux communiquer avant les fusions et à vérifier systématiquement les fonctions
importantes après leur résolution.

\subsection{Suivi des tâches}

Nous organisons le travail à l'aide de \textit{MeisterTask}. Chaque tâche possède un
responsable et un état d'avancement, ce qui permet d'identifier les fonctions à
intégrer et celles à vérifier. Les réunions hebdomadaires du jeudi matin servent à
faire le point, répartir les priorités et résoudre collectivement les blocages les
plus importants.

\begin{figure}[H]
    \centering
    \includegraphics[width=0.6\textwidth]{meistertask.png}
    \caption{Tableau de suivi des tâches sur MeisterTask}
\end{figure}

Dans les dernières phases, l'organisation des vérifications est devenue aussi importante que
l'ajout de fonctionnalités : une mécanique correcte à un joueur doit être revue à
deux, après une mort, après une reprise et après une transition de carte.

\sectionbreak
\section{Bilan et perspectives}

La version actuelle de \textit{Chamouraï} remplit le coeur de notre ambition : le
joueur apprend les contrôles dans une introduction guidée, explore un donjon peuplé
d'ennemis variés, récupère les clés auprès d'adversaires intermédiaires et termine sur un
combat en deux phases contre le golem. Le jeu est jouable seul ou à deux, dispose
d'une identité sonore et visuelle, et traite les principaux cas de victoire, de
défaite et de reprise.

Les évolutions les plus importantes ne sont pas seulement quantitatives. Le passage à
des cartes pilotant les apparitions et collisions facilite la création de contenu. La
synchronisation entre les deux joueurs permet à la progression de rester cohérente. Enfin, les
choix effectués pour le golem et l'apprentissage guidé montrent notre priorité : proposer une
expérience claire, équilibrée et agréable avant de multiplier les effets difficiles à
maîtriser.

Pour une version future, nous pourrions compléter les vérifications du réseau sur des connexions
moins stables et approfondir le combat contre le golem. Nous avions notamment imaginé
son attaque de rayon lumineux avant de l'abandonner, car les ressources graphiques
étaient difficiles à exploiter correctement avec Pygame ; l'onde de choc en arc de
cercle constitue la solution que nous avons su rendre fonctionnelle.

Nous aurions aussi aimé proposer plusieurs niveaux et permettre au joueur de
débloquer de nouvelles armes au fil de sa progression. Nous avons finalement préféré
concentrer nos efforts sur un unique niveau détaillé, avec une progression complète,
plusieurs adversaires et un combat final, plutôt que de multiplier des zones moins
abouties.

\imageplaceholder{Équipe ou écran final du projet}{Capture finale représentative du jeu terminé : victoire contre le golem, ou montage des principales scènes du niveau.}

\sectionbreak
\section{Annexes}

\begin{figure}[H]
    \centering
    \includegraphics[width=0.32\textwidth]{tuiles.png}
    \caption{Exemple de feuille de tuiles utilisée pour construire les cartes}
\end{figure}

\begin{figure}[H]
    \centering
    \includegraphics[width=0.22\textwidth]{instagram.png}
    \caption{Groupe Instagram de communication rapide}
\end{figure}

\imageplaceholder{Liste des ressources visuelles complémentaires}{Emplacement éventuel pour une planche réunissant les images des clés, du vieil homme, des adversaires majeurs et du golem.}

\end{document}
